\uuid{S3zU}
\exo7id{1598}
\titre{exo7 1598}
\auteur{legall}
\organisation{exo7}
\datecreate{1998-09-01}
\isIndication{false}
\isCorrection{false}
\chapitre{Réduction d'endomorphisme, polynôme annulateur}
\sousChapitre{Valeur propre, vecteur propre}
\cpge{mpsi-mp2i.calcul-matriciel-systemes.e, pcsi.determinants.c, ptsi.determinants.a, mpsi-mp2i.determinants.d, pcsi.calcul-matriciel-systemes-lineaires.d, ptsi.calcul-matriciel-systemes-lineaires.d, bcpst1.matrices, mp-mpi.endomorphismes-euclidien.b, pc.endomorphismes-euclidiens.b, psi.endomorphismes-euclidien.b, pt.isometries-euclidien.b, mp-mpi.reduction-endomorphismes.b, pc.reduction.a, psi.reduction-endomorphismes.a, pt.reduction-endomorphismes.a, bcpst2.valeurs-propres-vecteurs-propres.a}
\module{Algèbre}
\niveau{L2}
\difficulte{}

\contenu{
\question{
Soit $A \in \mathcal{O}_n(\R)$. Montrer que si $-1$ n'est pas valeur
propre de $A$, alors il existe une matrice $Q$ antisym\'etrique (i.e.
${}^tQ = -Q$) telle que $A = (I + Q)^{-1}(I - Q) = (I - Q)(I + Q)^{-1}$ et
qu'on a $A \in \mathcal{SO}_n(\R)$. R\'eciproque ?
}
}
